\section{Recognition of the smooth manifold}

\begin{lemma}\label{lem:homology-sphere}
Let $M$ be a closed simply connected $n$-manifold, $n\ge2$, with the
integral homology of $S^n$.  Then $M$ is homotopy equivalent to $S^n$.
\end{lemma}

\begin{proof}
If $k<n$ were the least integer with $\pi_k(M)\ne0$, the Hurewicz
theorem would give $\pi_k(M)\cong H_k(M;\Z)=0$, a contradiction.
Thus $\pi_i(M)=0$ for $1\le i<n$, and the Hurewicz theorem in degree
$n$ gives an isomorphism $\pi_n(M)\to H_n(M;\Z)\cong\Z$.  Choose
$u:S^n\to M$ representing a class mapping to a generator.  The map $u$
is an isomorphism on integral homology.  Since both spaces have the
homotopy type of CW complexes and are simply connected, the homology
Whitehead theorem implies that $u$ is a homotopy equivalence
\cite[Ch.~4]{Hatcher}.
\end{proof}

\begin{proof}[Proof of \cref{thm:main}]
By \cref{thm:pi1,thm:integral-cohomology,lem:homology-sphere}, the closed
smooth manifold underlying $X$ is a homotopy six-sphere.  For $n\ge5$, the group $\Theta_n$ consists of the oriented
$h$-cobordism classes of homotopy $n$-spheres.  In dimensions at least five, the smooth
$h$-cobordism theorem makes two such spheres diffeomorphic whenever they
represent the same class \cite{Milnor}.  Kervaire and Milnor computed
$\Theta_6=0$ \cite{KM}.  Hence $X$ is
orientation-preservingly diffeomorphic to the standard $S^6$.  Transporting
the complex atlas of $X$ through such a diffeomorphism gives an integrable
complex structure on $S^6$.
\end{proof}

